\section{Discussion: Evidence and Method Selection}
\label{ch:discussion}

This chapter separates what the experiments establish from what remains a
methodological hypothesis. The appropriate choice depends on the residual
dimension, estimator, bandwidth and bias, solver status, and independent
validation scope. A low held-out frequency is not a safety certificate when
the optimization gate has failed or when the fitted decision was selected from
the same data without a selection-aware bound.

\subsection{Answers to the research questions}
\label{sec:discussion-rq}

\subsubsection{RQ1: residual-space formulation}
\label{sec:discussion-rq1}

Chapters~3--5 establish the scalar residual reformulation and its analytic
gradient. For an individual event, estimating the pushforward residual CDF
avoids direct KDE in the uncertainty dimension. The claim is conditional on
the regularity and convergence assumptions stated in Chapter~5; it is not a
finite-sample population guarantee for the optimizer returned by the NLP.

\subsubsection{RQ2: numerical comparison}
\label{sec:discussion-rq2}

The static and dispatch experiments show a local objective--violation trade-off
among nominal, scenario, unbiased KDE, and shifted Epanechnikov formulations.
The dispatch adapter is the strongest transfer check because its mixture CDF is
known independently. The corrected training/test split in
Table~\ref{tab:dispatch-results} must be read together with solver status.
These are seeded synthetic comparisons, not universal rankings. Scenario
guarantees require convexity, IID sampling, support-rank and confidence-level
conditions that are not calculated by the local runner.

\subsubsection{RQ3: bandwidth, sample size, and bias}
\label{sec:discussion-rq3}

The paper-aligned records show how finite choices of $N$ and $h$ move the
returned decision. The biased rows have lower held-out violation and higher
objective in these runs, but the implementation also changes the selected
bandwidth and continuation branch. The observation therefore supports a local
safety--cost pattern, not a causal estimate of the bias alone and not a
convergence result.

\subsubsection{RQ4: joint-event representation}
\label{sec:discussion-rq4}

The maximum statistic represents the joint event exactly. Equal Boole and
maximum rows coincide in the current runner because they evaluate the same
Boolean event; this does not test componentwise risk allocation. The
log-sum-exp statistic satisfies $s_\tau\geq s$ and therefore changes the event
being measured. Its observed rate is a diagnostic of that changed statistic,
not evidence for a joint optimizer or a probability margin.

\subsection{Relation to the anchor literature}
\label{sec:discussion-literature}

\subsubsection{What the Gaussian sequence tests}
\label{sec:discussion-schuster}

The local nested Gaussian sequence probes the direction of the Schuster
convergence hypothesis. It does not compute uniform residual-grid error,
verify the theorem's geometric assumptions, or establish a limiting optimizer.
The dispatch adapter is the closest local analogue to the convex setting; the
lunar transcription is not.

\subsubsection{What the biased arm does not replicate}
\label{sec:discussion-keil}

The local shifted Epanechnikov arm tests the one-sided kernel property with IID
samples and finite-dimensional SLSQP. The lunar adapter additionally uses an
Euler mesh, whereas the cited Keil workflow uses MCMC, LGR, and continuation.
The local arm is therefore an empirical analogue, not a replication of that
workflow.

\subsection{Mechanistic synthesis}
\label{sec:discussion-mechanisms}

\subsubsection{Statistical mechanism}
\label{sec:mechanism-statistical}

Symmetric smoothing can credit positive residuals near the boundary. The
$B(h)=h$ Epanechnikov shift removes this per-sample optimism for its compact
support, but the resulting inequality is an expectation statement for a fixed,
sample-independent decision. Calibration or a selection-aware concentration
bound is still needed after optimization.

\subsubsection{Optimization mechanism}
\label{sec:mechanism-optimization}

The active estimated constraint and the objective determine the returned
decision together. The observed biased premium is therefore a property of the
full estimator--bandwidth--continuation pipeline. A controlled ablation must
hold $h$ and the continuation schedule fixed before attributing an objective or
violation change to the kernel bias.

\subsubsection{Transfer and solver mechanism}
\label{sec:mechanism-transfer}

The exact dispatch CDF provides a distributional oracle for the affine balance
residual. In contrast, the robust lunar stages have failed solver status and
terminal-only validation; the Euler/trapezoidal mismatch also changes the
stored objective. Those rows are valuable audit records but cannot support a
physical cost or continuous-time safety claim.

\subsubsection{Joint-event mechanism}
\label{sec:mechanism-joint}

The maximum and equal-Boole representations coincide only because the current
runner evaluates the same event. The LSE upper bound smooths the max while
introducing a different scalar event. A genuine allocation study must retain
component residuals, optimize separate component constraints, and validate the
union event independently.

\subsection{Method-selection boundaries}
\label{sec:discussion-selection}

\subsubsection{Scenario and parametric alternatives}
\label{sec:selection-alternatives}

Scenario optimization is appropriate when its convexity, IID, support-rank, and
confidence assumptions are checked. A parametric reformulation is appropriate
when diagnostics support the assumed law and a deterministic equivalent is
available.

\subsubsection{When KDE is defensible}
\label{sec:selection-kde}

Unbiased KDE is defensible for low-dimensional empirical residuals when smooth
derivatives are valuable and bandwidth, solver margins, and independent test
performance are reported. One-sided compact kernels are defensible when the
fixed-decision envelope and its objective premium are worth the additional
conservatism. High-dimensional, rare-event, dependent, or shifted settings
require distributionally robust, importance-sampling, or tail-specific methods
before direct KDE is used.

\subsection{Future Work}\label{sec:discussion-future-work}


The results identify a focused doctoral research programme. Each item below is a
research target rather than a result established by this thesis.

\subsubsection{Selection-aware finite-sample reliability}

Chapter~4 gives an expectation envelope for a fixed decision. The principal next
step is to make a corresponding statement for the decision returned after the
same data have been used to fit the KDE and solve the nonlinear programme. An
initial, conservative route is to separate training, calibration, and test data:
the optimizer uses the training sample, while an independent calibration sample
sets a lower confidence bound on the achieved safe probability. A sharper route
would prove a uniform bound over the decision class and a prespecified grid of
bandwidths, temperatures, and component weights. The bound should separate
sampling error, one-sided kernel bias, and the log-sum-exp approximation, and it
should state explicitly how its complexity grows with the number $m$ of joint
constraints. Such a result would turn the current fixed-decision envelope into a
post-selection safety statement with confidence level $1-\delta$.

The proposed target can be stated using the scalar statistic $s$ and its
log-sum-exp approximation $s_\tau$ from Chapter~4. For a training sample
$\xi_1,\ldots,\xi_N$, write the corresponding biased estimator as
\begin{equation}
  \widehat p^{\,b}_{N,h,\tau}(x)
  = \frac{1}{N}\sum_{j=1}^{N}
    F_K\!\left(
      \frac{-s_\tau(x,\xi_j)-b(h)}{h}
    \right),
  \qquad
  p(x)=\Prob\{s(x,\xi)\leq 0\} .
  \label{eq:future-postselection-estimator}
\end{equation}
The theorem to seek is a post-selection implication, uniform over $x$ and over
a prespecified tuning set $\mathcal H\times\mathcal T$:
\begin{equation}
  \widehat p^{\,b}_{N,h,\tau}(\widehat x)
  \geq 1-\varepsilon+\Delta_{N,h,\tau}(\delta)
  \quad\Longrightarrow\quad
  p(\widehat x)\geq1-\varepsilon
  \qquad\text{with probability at least }1-\delta .
  \label{eq:future-postselection-bound}
\end{equation}
Here the required margin is a research object, not a quantity established in
this thesis. A possible decomposition is only a conjectural template: the
$\tau\log m$ term is a residual-statistic margin and becomes a probability
margin only under an anti-concentration condition. The constants, exponent,
and complexity term must be derived rather than assumed:
\begin{equation}
  \Delta_{N,h,\tau}(\delta)
  = \underbrace{C_{\rm stat}
    \sqrt{\frac{\mathcal C(\mathcal H,\mathcal T,X)+\log(1/\delta)}{N}}}
    _{\text{sampling and selection}}
  + \underbrace{C_{\rm kde}h^{\beta}}
    _{\text{kernel smoothing}}
  + \underbrace{C_{\rm lse}\tau\log m}
    _{\text{joint smoothing}} ,
  \label{eq:future-margin-decomposition}
\end{equation}
where the constants and the complexity term would have to be derived from
explicit residual regularity, anti-concentration, and function-class assumptions.
An independent calibration theorem is a rigorous first milestone; a same-sample
result with a non-vacuous complexity term is the stronger doctoral target.

\subsubsection{Conservatism and objective cost}

Safety is useful only if its objective cost is quantified. Under a density
regularity or anti-concentration condition near the residual boundary, future
work should bound the probability mass introduced by bandwidth $h$ and
temperature $\tau$, then propagate that bound through an objective error bound
or sharp-minimum condition. This would provide a principled rule for choosing
$(h,\tau)$ and would establish when the biased construction approaches the true
optimum as $N$ increases. The claim must remain conditional: uniqueness alone is
insufficient without the compactness, coercivity, or error-bound assumptions used
in the optimization argument.

With an appropriate objective error bound, the desired conclusion would take the
form
\begin{equation}
  0\leq f(\widehat x_{N,h,\tau})-f(x^*)
  \leq C_{\rm opt}\,\Delta_{N,h,\tau}(\delta),
  \label{eq:future-objective-gap}
\end{equation}
for a solution $x^*$ of the exact problem. Establishing this inequality, rather
than merely observing a higher objective in a finite experiment, would quantify
the price of one-sided protection.

\subsubsection{Continuous-time chance constraints}

The lunar experiment currently evaluates a discretized path. A stronger control
result would repair the transcription and quadrature first, then derive a
mesh-to-trajectory residual bound from temporal regularity and integrator error.
The statistical safety margin could subsequently be combined with that
discretization margin to control the probability of a violation at any time in
the horizon. This requires stating the dependence structure of the disturbance
trajectory and validating the complete path on a refined mesh; nodewise test
frequencies cannot establish this claim.

\subsubsection{Validation and transfer}

The theoretical targets should be evaluated with independent repetitions rather
than a single seed. Static and dispatch problems with exact probability oracles
should report coverage, excess violation, objective regret, solver status, and
runtime over a factorial grid of $N$, $h$, $\tau$, and $m$. The joint experiment
should optimize the componentwise event itself, rather than only compare scalar
statistics. Finally, dependence, distribution shift, and a public or
semi-real benchmark are needed to determine whether the observed conservatism
survives outside the present synthetic setting.
